% GENERATED FILE. Edit canonical lessons, then run npm run generate:books.
% canonical-source-hash: fnv1a64:b0a2732858397d97
% canonical-lessons: KA-C418-paryaya, KA-C418-ayke, KA-C418-tatkalika, KA-C418-sasvata, KA-C418-marpadisu, KA-R418-review-412-418

\chapter{Other Options}
\label{ch:ka-other-options}
\hlchaptermodality{418}

\begin{chapteropening}
\textbf{By the end of this chapter:} \emph{I can say an alternative, a choice, temporary, permanent and to modify.}

Complete the last lesson of chapter 418: I can say an alternative, a choice, temporary, permanent and to modify.
\end{chapteropening}

\section[\texorpdfstring{\kn{ಪರ್ಯಾಯ} (paryāya)}{ಪರ್ಯಾಯ (paryāya)}]{\kn{ಪರ್ಯಾಯ} (paryāya) — an alternative}

\label{lesson:KA-C418-paryaya}



\begin{quote}
Before the new one: say the Kannada for to recover dues, then the Kannada for to submit.
\end{quote}

\subsection*{You'll want to know: \kn{ಪರ್ಯಾಯ}}
\textbf{\kn{ಪರ್ಯಾಯ}} — \emph{paryāya} — \textquotedblleft{}an alternative\textquotedblright{}.

\begin{quote}
\kn{ಬಸ್} \kn{ಇಲ್ಲದಿದ್ದರೆ} \kn{ಬೇರೆ} \kn{ಪರ್ಯಾಯ} \kn{ಇಲ್ಲ}.
\end{quote}

\begin{quote}
\emph{bas illadiddare bēre paryāya illa.}
\end{quote}

\begin{quote}
If there is no bus, there is no other option.
\end{quote}

\kn{ಪರ್ಯಾಯವಾಗಿ}, as an alternative.

\begin{cousinweb}
\textbf{\kn{ಪರ್ಯಾಯ}} (\emph{paryāya}) is borrowed from Sanskrit \emph{paryāya}, \textquotedblleft{}going round, succession, turn\textquotedblright{} (Monier-Williams).
\end{cousinweb}

\begin{grammarlens}[title={What you've built}]
One more word for saying what happens if.
\end{grammarlens}

\subsection*{Your turn}
\begin{itemize}
\raggedright
  \item \emph{Say it:} \emph{paryāya}
  \item \emph{Say it:} \emph{paryāya}, once more
  \item \emph{Say it:} \emph{bas illadiddare bēre paryāya illa.} — If there is no bus, there is no other option.
  \item \emph{Recall:} say the Kannada for to recover dues, then the Kannada for to submit, then say \emph{paryāya} again
\end{itemize}

\begin{tcolorbox}[breakable,colback=teal!4,colframe=teal!35!black,title={Before you move on}]
What does \kn{ಪರ್ಯಾಯ} mean? (\textquotedblleft{}An alternative\textquotedblright{}.) Say it once more.
\end{tcolorbox}
\section[\texorpdfstring{\kn{ಆಯ್ಕೆ} (āyke)}{ಆಯ್ಕೆ (āyke)}]{\kn{ಆಯ್ಕೆ} (āyke) — a choice; a selection}

\label{lesson:KA-C418-ayke}



\begin{quote}
Before the new one: say the Kannada for to submit, then the Kannada for an alternative.
\end{quote}

\subsection*{You'll want to know: \kn{ಆಯ್ಕೆ}}
\textbf{\kn{ಆಯ್ಕೆ}} — \emph{āyke} — \textquotedblleft{}a choice; a selection\textquotedblright{}.

\begin{quote}
\kn{ಯಾವ} \kn{ಬಣ್ಣ} \kn{ಬೇಕು} \kn{ಅನ್ನುವುದು} \kn{ನಿನ್ನ} \kn{ಆಯ್ಕೆ}.
\end{quote}

\begin{quote}
\emph{yāva baṇṇa bēku annuvudu ninna āyke.}
\end{quote}

\begin{quote}
Which colour you want is your choice.
\end{quote}

\kn{ಆಯ್ಕೆ} \kn{ಮಾಡು}, to choose.

\begin{cousinweb}
\textbf{\kn{ಆಯ್ಕೆ}} (\emph{āyke}) is native Dravidian, a noun from \emph{āy}, \textquotedblleft{}to collect, to select\textquotedblright{} (DEDR 363).
\end{cousinweb}

\begin{grammarlens}[title={What you've built}]
One more word for saying what happens if.
\end{grammarlens}

\subsection*{Your turn}
\begin{itemize}
\raggedright
  \item \emph{Say it:} \emph{āyke}
  \item \emph{Say it:} \emph{āyke}, once more
  \item \emph{Say it:} \emph{yāva baṇṇa bēku annuvudu ninna āyke.} — Which colour you want is your choice.
  \item \emph{Recall:} say the Kannada for to submit, then the Kannada for an alternative, then say \emph{āyke} again
\end{itemize}

\begin{tcolorbox}[breakable,colback=teal!4,colframe=teal!35!black,title={Before you move on}]
What does \kn{ಆಯ್ಕೆ} mean? (\textquotedblleft{}A choice; a selection\textquotedblright{}.) Say it once more.
\end{tcolorbox}
\section[\texorpdfstring{\kn{ತಾತ್ಕಾಲಿಕ} (tātkālika)}{ತಾತ್ಕಾಲಿಕ (tātkālika)}]{\kn{ತಾತ್ಕಾಲಿಕ} (tātkālika) — temporary}

\label{lesson:KA-C418-tatkalika}



\begin{quote}
Before the new one: say the Kannada for an alternative, then the Kannada for a choice.
\end{quote}

\subsection*{You'll want to know: \kn{ತಾತ್ಕಾಲಿಕ}}
\textbf{\kn{ತಾತ್ಕಾಲಿಕ}} — \emph{tātkālika} — \textquotedblleft{}temporary\textquotedblright{}.

\begin{quote}
\kn{ಇದು} \kn{ತಾತ್ಕಾಲಿಕ} \kn{ಕೆಲಸ}, \kn{ಮೂರು} \kn{ತಿಂಗಳು} \kn{ಮಾತ್ರ}.
\end{quote}

\begin{quote}
\emph{idu tātkālika kelasa, mūru tiṅgaḷu mātra.}
\end{quote}

\begin{quote}
This is temporary work, only three months.
\end{quote}

\kn{ತಾತ್ಕಾಲಿಕವಾಗಿ}, for the time being.

\begin{cousinweb}
\textbf{\kn{ತಾತ್ಕಾಲಿಕ}} (\emph{tātkālika}) is borrowed from Sanskrit \emph{tātkālika}, \textquotedblleft{}lasting that time\textquotedblright{} (Monier-Williams), from \emph{tat-kāla}, \textquotedblleft{}that time\textquotedblright{}.
\end{cousinweb}

\begin{grammarlens}[title={What you've built}]
One more word for saying what happens if.
\end{grammarlens}

\subsection*{Your turn}
\begin{itemize}
\raggedright
  \item \emph{Say it:} \emph{tātkālika}
  \item \emph{Say it:} \emph{tātkālika}, once more
  \item \emph{Say it:} \emph{idu tātkālika kelasa, mūru tiṅgaḷu mātra.} — This is temporary work, only three months.
  \item \emph{Recall:} say the Kannada for an alternative, then the Kannada for a choice, then say \emph{tātkālika} again
\end{itemize}

\begin{tcolorbox}[breakable,colback=teal!4,colframe=teal!35!black,title={Before you move on}]
What does \kn{ತಾತ್ಕಾಲಿಕ} mean? (\textquotedblleft{}Temporary\textquotedblright{}.) Say it once more.
\end{tcolorbox}
\section[\texorpdfstring{\kn{ಶಾಶ್ವತ} (śāśvata)}{ಶಾಶ್ವತ (śāśvata)}]{\kn{ಶಾಶ್ವತ} (śāśvata) — permanent}

\label{lesson:KA-C418-sasvata}



\begin{quote}
Before the new one: say the Kannada for a choice, then the Kannada for temporary.
\end{quote}

\subsection*{You'll want to know: \kn{ಶಾಶ್ವತ}}
\textbf{\kn{ಶಾಶ್ವತ}} — \emph{śāśvata} — \textquotedblleft{}permanent\textquotedblright{}.

\begin{quote}
\kn{ಅವನಿಗೆ} \kn{ಶಾಶ್ವತ} \kn{ಕೆಲಸ} \kn{ಸಿಕ್ಕಿತು}.
\end{quote}

\begin{quote}
\emph{avanige śāśvata kelasa sikkitu.}
\end{quote}

\begin{quote}
He got a permanent job.
\end{quote}

\kn{ಶಾಶ್ವತವಾಗಿ}, for good.

\begin{cousinweb}
\textbf{\kn{ಶಾಶ್ವತ}} (\emph{śāśvata}) is borrowed from Sanskrit \emph{śāśvata}, \textquotedblleft{}eternal, constant\textquotedblright{} (Monier-Williams).
\end{cousinweb}

\begin{grammarlens}[title={What you've built}]
One more word for saying what happens if.
\end{grammarlens}

\subsection*{Your turn}
\begin{itemize}
\raggedright
  \item \emph{Say it:} \emph{śāśvata}
  \item \emph{Say it:} \emph{śāśvata}, once more
  \item \emph{Say it:} \emph{avanige śāśvata kelasa sikkitu.} — He got a permanent job.
  \item \emph{Recall:} say the Kannada for a choice, then the Kannada for temporary, then say \emph{śāśvata} again
\end{itemize}

\begin{tcolorbox}[breakable,colback=teal!4,colframe=teal!35!black,title={Before you move on}]
What does \kn{ಶಾಶ್ವತ} mean? (\textquotedblleft{}Permanent\textquotedblright{}.) Say it once more.
\end{tcolorbox}
\section[\texorpdfstring{\kn{ಮಾರ್ಪಡಿಸು} (mārpaḍisu)}{ಮಾರ್ಪಡಿಸು (mārpaḍisu)}]{\kn{ಮಾರ್ಪಡಿಸು} (mārpaḍisu) — to modify}

\label{lesson:KA-C418-marpadisu}



\begin{quote}
Before the new one: say the Kannada for temporary, then the Kannada for permanent.
\end{quote}

\subsection*{You'll want to know: \kn{ಮಾರ್ಪಡಿಸು}}
\textbf{\kn{ಮಾರ್ಪಡಿಸು}} — \emph{mārpaḍisu} — \textquotedblleft{}to modify\textquotedblright{}.

\begin{quote}
\kn{ದರ್ಜಿ} \kn{ನನ್ನ} \kn{ಅಂಗಿಯನ್ನು} \kn{ಸ್ವಲ್ಪ} \kn{ಮಾರ್ಪಡಿಸಿದನು}.
\end{quote}

\begin{quote}
\emph{darji nanna aṅgiyannu svalpa mārpaḍisidanu.}
\end{quote}

\begin{quote}
The tailor altered my shirt a little.
\end{quote}

Past (I …): \kn{ಮಾರ್ಪಡಿಸಿದೆ}, \emph{mārpaḍiside}.

\kn{ಮಾರ್ಪಾಡು}, a change.

\begin{cousinweb}
\textbf{\kn{ಮಾರ್ಪಡಿಸು}} (\emph{mārpaḍisu}) is the causative of \emph{mārpaḍu}, \textquotedblleft{}to change\textquotedblright{}, and is usually traced to native \emph{māṟu}, \textquotedblleft{}to be opposite, to take or give one thing for another\textquotedblright{} (DEDR 4834), the root of \kn{ಮಾರು}, \textquotedblleft{}to sell\textquotedblright{}; DEDR does not list the verb itself, so the formation is not settled.
\end{cousinweb}

\begin{grammarlens}[title={What you've built}]
One more word for saying what happens if.
\end{grammarlens}

\subsection*{Your turn}
\begin{itemize}
\raggedright
  \item \emph{Say it:} \emph{mārpaḍisu}
  \item \emph{Say it:} \emph{mārpaḍisu}, once more
  \item \emph{Say it:} \emph{darji nanna aṅgiyannu svalpa mārpaḍisidanu.} — The tailor altered my shirt a little.
  \item \emph{Recall:} say the Kannada for temporary, then the Kannada for permanent, then say \emph{mārpaḍisu} again
\end{itemize}

\begin{tcolorbox}[breakable,colback=teal!4,colframe=teal!35!black,title={Before you move on}]
What does \kn{ಮಾರ್ಪಡಿಸು} mean? (\textquotedblleft{}To modify\textquotedblright{}.) Say it once more.
\end{tcolorbox}
\section[(dialogue)]{Looking back: from When Danger Comes to Other Options}

\label{lesson:KA-R418-review-412-418}



\begin{quote}
Say the Kannada for a danger and a precaution.
\end{quote}

\subsection*{Your turn}
Say these aloud:

\begin{itemize}
\raggedright
  \item a danger, a precaution, probable, hardship, to ward off
  \item a lottery, gambling, a bet, to estimate, an average
  \item strictness, maximum, minimum, to fix, to restrict
  \item suspension, to exclude, an exemption, to sanction, an official order
  \item a written undertaking, a vow, a reward, to ask formally, an oath of office
  \item an instalment, arrears, savings, to recover dues, to submit
  \item an alternative, a choice, temporary, permanent, to modify
\end{itemize}

\begin{tcolorbox}[breakable,colback=teal!4,colframe=teal!35!black,title={Before you move on}]
A danger? (\textbf{\kn{ಅಪಾಯ}}, \emph{apāya}.) A precaution? (\textbf{\kn{ಮುನ್ನೆಚ್ಚರಿಕೆ}}, \emph{munneccarike}.) Probable? (\textbf{\kn{ಸಂಭವನೀಯ}}, \emph{sambhavanīya}.) Hardship? (\textbf{\kn{ಸಂಕಷ್ಟ}}, \emph{saṅkaṣṭa}.) To ward off? (\textbf{\kn{ನಿವಾರಿಸು}}, \emph{nivārisu}.) A lottery? (\textbf{\kn{ಲಾಟರಿ}}, \emph{lāṭari}.) Gambling? (\textbf{\kn{ಜೂಜು}}, \emph{jūju}.) A bet? (\textbf{\kn{ಪಣ}}, \emph{paṇa}.) To estimate? (\textbf{\kn{ಅಂದಾಜಿಸು}}, \emph{andājisu}.) An average? (\textbf{\kn{ಸರಾಸರಿ}}, \emph{sarāsari}.) Strictness? (\textbf{\kn{ಕಟ್ಟುನಿಟ್ಟು}}, \emph{kaṭṭuniṭṭu}.) Maximum? (\textbf{\kn{ಗರಿಷ್ಠ}}, \emph{gariṣṭha}.) Minimum? (\textbf{\kn{ಕನಿಷ್ಠ}}, \emph{kaniṣṭha}.) To fix? (\textbf{\kn{ನಿಗದಿಪಡಿಸು}}, \emph{nigadipaḍisu}.) To restrict? (\textbf{\kn{ನಿರ್ಬಂಧಿಸು}}, \emph{nirbandhisu}.) Suspension? (\textbf{\kn{ಅಮಾನತು}}, \emph{amānatu}.) To exclude? (\textbf{\kn{ಹೊರಗಿಡು}}, \emph{horagiḍu}.) An exemption? (\textbf{\kn{ವಿನಾಯಿತಿ}}, \emph{vināyiti}.) To sanction? (\textbf{\kn{ಮಂಜೂರು} \kn{ಮಾಡು}}, \emph{mañjūru māḍu}.) An official order? (\textbf{\kn{ಆದೇಶ}}, \emph{ādēśa}.) A written undertaking? (\textbf{\kn{ಮುಚ್ಚಳಿಕೆ}}, \emph{muccaḷike}.) A vow? (\textbf{\kn{ಪ್ರತಿಜ್ಞೆ}}, \emph{pratijñe}.) A reward? (\textbf{\kn{ಪ್ರತಿಫಲ}}, \emph{pratiphala}.) To ask formally? (\textbf{\kn{ಕೋರು}}, \emph{kōru}.) An oath of office? (\textbf{\kn{ಪ್ರಮಾಣವಚನ}}, \emph{pramāṇavacana}.) An instalment? (\textbf{\kn{ಕಂತು}}, \emph{kantu}.) Arrears? (\textbf{\kn{ಬಾಕಿ}}, \emph{bāki}.) Savings? (\textbf{\kn{ಉಳಿತಾಯ}}, \emph{uḷitāya}.) To recover dues? (\textbf{\kn{ವಸೂಲಿ} \kn{ಮಾಡು}}, \emph{vasūli māḍu}.) To submit? (\textbf{\kn{ಸಲ್ಲಿಸು}}, \emph{sallisu}.) An alternative? (\textbf{\kn{ಪರ್ಯಾಯ}}, \emph{paryāya}.) A choice? (\textbf{\kn{ಆಯ್ಕೆ}}, \emph{āyke}.) Temporary? (\textbf{\kn{ತಾತ್ಕಾಲಿಕ}}, \emph{tātkālika}.) Permanent? (\textbf{\kn{ಶಾಶ್ವತ}}, \emph{śāśvata}.) To modify? (\textbf{\kn{ಮಾರ್ಪಡಿಸು}}, \emph{mārpaḍisu}.)
\end{tcolorbox}
